\section{Thagaste: A Family on the Edge of Two Worlds}

\subsection{Birth and household}

Aurelius Augustinus was born on 13 November 354 at Thagaste in the Roman
province of Numidia---the modern Souk Ahras in Algeria---to Patricius, a
pagan town councillor of modest means who became a catechumen late in life,
and Monica, a Christian from a Christian family (Benedict XVI, General
Audience, 9 January 2008; \work{Confessions} 2.3.5--8, 9.9.19--22). Augustine
himself calls his father ``but a poor freeman of Thagaste'' whose ambition
outran his purse: the funds for Carthage were found ``rather by the
determination than the means of my father,'' whom the neighbors praised for
going ``even beyond his means to supply his son with all the necessaries for
a far journey for the sake of his studies.'' Augustine records the sacrifice
and its limit in one sentence: ``But yet this same father did not trouble
himself how I grew towards You, nor how chaste I was, so long as I was
skilful in speaking'' (\work{Confessions} 2.3.5). The household
spoke Latin; Augustine mastered it completely, learned Greek reluctantly and
imperfectly, and never learned the Punic of the countryside (Benedict XVI,
General Audience, 9 January 2008).

As an infant he was enrolled, in the African manner, as a candidate for
baptism rather than baptized: ``I was signed with the sign of the cross, and
was seasoned with His salt even from the womb of my mother, who greatly
trusted in You.'' When as a boy he fell dangerously ill and begged for
baptism, the rite was deferred lest post-baptismal sin prove worse: ``So my
cleansing was deferred, as if I must needs, should I live, be further
polluted'' (\work{Confessions} 1.11.17). The deferral was normal practice;
Augustine the bishop recorded it with disapproval, and it set the shape of
his story---a Christian catechumen by cradle, a baptized Christian only at
thirty-two.

\subsection{The pear tree}

Of the schooling at Thagaste and nearby Madauros the \work{Confessions}
preserve the memories a penitent selects: beatings, the seduction of pagan
literature, weeping for Dido while ignoring his own soul. One episode from
his sixteenth year, an idle year at home while his father saved for
Carthage, became the most famous minor theft in literature. ``I pilfered
that of which I had already sufficient, and much better. Nor did I desire to
enjoy what I pilfered, but the theft and sin itself. There was a pear-tree
close to our vineyard, heavily laden with fruit, which was tempting neither
for its colour nor its flavour''; he and his companions shook it down late
at night and ``carried away great loads, not to eat ourselves, but to fling
to the very swine, having only eaten some of them; and to do this pleased us
all the more because it was not permitted.'' The grown man circles the
memory for chapters, appalled at ``having no inducement to evil but the evil
itself'' (\work{Confessions} 2.4.9). The same
book preserves the domestic texture of the year: Patricius seeing his son at
the baths and rejoicing at the prospect of grandchildren; Monica warning him
``not to commit fornication; but above all things never to defile another
man's wife'' (\work{Confessions} 2.3.6--8).

\witnesslimit{The pear-tree episode is subject-authored evidence of class A:
an event of c.~370 narrated c.~397--400 as a theological set piece on the
gratuity of sin. Its factual core---a boys' theft of unwanted fruit---is
unremarkable and uncontrolled by any other witness; its enormous later
fame belongs to the history of the \work{Confessions}' reception, not to
the event. This study reports it as Augustine's own confession, told for
the purpose he states.}

\subsection{Carthage: the cauldron and the concubine}

At about seventeen, funded by his father---whom Monica, ``towards the end
of his earthly existence, did\ldots\ gain over unto You,'' and who died soon
after (\work{Confessions} 9.9.22)---and by the patronage of the local
magnate Romanianus, Augustine went to Carthage to study rhetoric. His own retrospect opens with the
sentence every biographer since has quoted: ``To Carthage I came, where a
cauldron of unholy loves bubbled up all around me'' (\work{Confessions}
3.1.1). The reality the same book documents is more settled than the
sentence. Within a year or so he had taken a concubine, and he kept faith
with her for some fourteen years: ``In those years I had one (whom I knew
not in what is called lawful wedlock, but whom my wayward passion, void of
understanding, had discovered), yet one only, remaining faithful even to
her'' (\work{Confessions} 4.2.2). About 372 she bore him a son, Adeodatus
---``the son of my sin,'' as he later wrote, and also, in the same breath,
one of God's good gifts (\work{Confessions} 9.6.14). Concubinage of this
kind was a recognized Roman institution for a rising man who expected one
day to marry property; Augustine's culpability, in his own later judgment,
lay not in scandal but in the divided will the arrangement embodied.

\witnesslimit{The woman's name is unknown. Augustine, the only witness,
never gives it---not at the union (\work{Confessions} 4.2.2), not at the
separation (6.15.25)---and no other ancient source supplies it. Every name
in later fiction is invention. The silence is itself evidence, though of
what---discretion, shame, reverence, or rhetorical design---cannot be
determined, and this study asserts nothing about her beyond what the
\work{Confessions} state.}

\subsection{The \work{Hortensius}}

In his nineteenth year, ``the ordinary course of study'' brought him to a
now-lost dialogue of Cicero exhorting the reader to philosophy. ``This book,
in truth, changed my affections, and turned my prayers to Yourself, O Lord,
and made me have other hopes and desires\ldots\ with an incredible warmth of
heart, I yearned for an immortality of wisdom'' (\work{Confessions} 3.4.7).
He turned to the Scriptures and bounced off them: to a connoisseur of
Cicero ``they appeared to me to be unworthy to be compared with the dignity
of Tully'' (\work{Confessions} 3.5.9). The conversion to
philosophy without a Church prepared the next step, and Augustine took it
promptly: ``Therefore I fell among men proudly raving, very carnal, and
voluble, in whose mouths were the snares of the devil''---the Manichees
(\work{Confessions} 3.6.10).
